\section{Conclusion}
\label{sec:conclusion}

The selected inverse is a classical, exact, $\bigO(\text{fill})$ global operator, and,
because it is the gradient of a log-determinant, it is cheap to
differentiate. Those facts are not new, and this paper's own numerics rule
out the obvious payoff: exactness buys nothing on the precision axis
(Section~\ref{sec:numerics}). The hypothesis under test was that it buys
something else instead -- that exactness is an inductive bias, changing which
of two equally-fit models is preferred. Two pre-registered, matched-fit
probes on $30$ held-out seeds test this directly, and on both, the
inductive-bias branch does not fire: the alternative (bounded-reach sweeps in
Section~\ref{sec:maze}, the biased Neumann adjoint in Section~\ref{sec:deq})
cannot reach comparable training fit at the difficulty levels where the raw
gap is largest, so the gap is explained by attainability, not by an implicit
preference among optima that fit equally well. This is a real finding, not a
null one: it says the practical advantage of embedding the exact solve --
still a factor of roughly $600\times$ to $28{,}000\times$ across the tested
sizes -- comes from
the alternatives' inability to represent or optimize the target, not from a
training-time bias effect. Both probes' negative controls are disclosed as
imperfect (Section~\ref{sec:program}), so even this conservative reading is
reported as hedged, following the frozen pre-registration rather than
relaxing it post hoc.

The limits are explicit: the evidence is at controlled CPU scale and
confined to low-treewidth (1-D and quasi-2-D) problems. What remains open is
whether exactness matters for reasons unrelated to inductive bias: exact
deterministic gradients for calibrated uncertainty quantification,
conditioning-independent adjoints for differentiable physics, and a
one-layer exact resolvent against message-passing's bounded-depth limits
(Section~\ref{sec:program}). These are numerical and expressivity claims,
each with its own falsification bar, gated on a single scoped build: a
GPU-batched kernel. The next step is to run them. If they confirm the claim,
exactness's numerical properties become a design principle for
solve-shaped learning at scale; if they fail, they pin down where exactness
does not matter, which is itself an answer. What this paper does not
support, on its own pre-registered evidence, is the inductive-bias framing
that motivated it.
